\section{Introduction}

\subsection{Background and Motivation}
The proliferation of autonomous mobile systems has revolutionized sectors
ranging from industrial logistics to domestic services. At the core of these 
systems lies visual perception and control, which enables robots to perceive and 
interpret unstructured environments and execute complex navigation tasks[1]. In recent
years, Convolutional Neural Networks(CNNs) have established themselves as the state-of-art solution
for robotic vision due to their superior performance in feature extraction 
and pattern recognition[2].\\

However, the deployment of Deep Learning models on some mobile robotic platforms
(most of them are resource-constrained) introduces a fundamental trade-off between 
computational performance and energy efficiency. Unlike cloud-based AI systems with
almost unlimited computational and energy resources, mobile robots and 
embedded systems-- especially micro-robots and self-balancing robots-- 
operate under significant constraints and challenges due to their limites computational resources,
including memory, processing power, and energy consumption[3].

The continuous processing of 
computionally intensive CNNs on battery-powered embedded processors often leads to
rapid depletion and thermal throttling, which severely limites the operational duration
of the robot and its ability to perform real-time tasks.\\

To address these challenges, Neuromorphic Computing has emerged as a promising paradigm
shift in embedded Artificial Intelligence. Inspired by the energy-efficient mechanisms
of the biological brain, Spiking Neural Networks(SNNs) process information using discrete,
sparse events(Spikes) rather than continuous numerical values, can offer
order-of-magnitude improments in energy efficiency compared to traditional Artificial Neural
Networks(ANNs)[5]. Consequently, investigating the application of SNNs in Edge Robotics
represents a critical frontier in modern robotic engineering.

\subsection{Problem Statement}
Despite their theoretical potential, the pratical deployment of SNNs on 
physical rototic platforms faces several critical engineering challenges:
\begin{description}
    \item[Computational Efficiency on Non-Neuromorphic Hardware:] While SNNs are 
    designed for specialized neuromorphic 
    chips (e.g., Intel Loihi), most edge robots, such as the Yahboom self-balancing 
    platform,
     rely on general-purpose microcontrollers like the STM32 and ESP32.
      Achieving high-performance SNN inference on these platforms requires rigorous 
      optimization of software-hardware co-design.
    
    \item[Sensory Integration and Real-time Control:] Autonomous systems such as two-wheeled
    self-balancing robots are dynamically unstable and require millisecond-level feedback loops. 
    Conventional frame-based vision sensors generate redundant data, which contradicts 
    the event-driven nature of SNNs. Developing a visual encoding mechanism that can bridge 
    frame-based input with spike-based processing isessential for maintaining system stability.

    
    \item[The Sim-to-Real Gap:] Many SNN studies remain confined to theoretical simulations
    or static dataset benchmarks. There is a pressing need for empirical research that demonstrates
    the feasibility of SNNs in managing closed-loop, real-time control tasks on physical hardware.
\end{description}

\subsection{Research Aims and Objectives}
The primary aim of this research is to design and implement a high-performance,
energy-efficient visual control framework for a self-balancing robot, leveraging the sparse compurting
characteristics of SNNs. The research will focus on achieving real-time, colsed-loop
navigation by optimizing the synergy between neuromorphic algorithms and embedded hardware.
To achieve this aim, the following objectives have been established:
\begin{enumerate}[label=\textbf{Objective \arabic*:}, leftmargin=*]
    \item \textbf{Development of a Hierarchical Control Architecture}\newline
    To establish a multi-layer control system using ROS 2 Humble. This involves
    utilizing an ESP32-S3 as a high-speed communication gateway to bridge the high-level
    SNN cognitive logic(running on a host machine) with the low-level reflexive control
    (implemented on an STM32F103). The architecture will leverage micro-ROS to ensure 
    synchronized data exchange between the SNN decision-making and the PID balancing loops.
    
    
    \item \textbf{Design of a Retina-inspired Visual Encoder}\newline
    To develop a visual encoding scheme based on temporal difference algorithms. 
    By computing the pixel-level intensity changes between successive frames,
    the system converts standaed RGB video streams into sparse, event-based spike
    tensors. This process aims to minimize redundant computation at the edge, mirroring
    the biological efficiency of the human retina.
    
    \item \textbf{Implementation of an Optimized C++ Inference Engine}\newline
    To engineer a dedicated SNN inference engine using the Eigen library within a C++
    environment. Utilizing software-hardware co-design techniques, this objective 
    focuses on minimizing processing latency to meet the millisecond-level responsiveness
    required for dynamically unstable platforms.

    \item \textbf{Performance Evalution and Benchmarking}\newline
    To quantitatively access the navigation accuracy and energy efficiency of the system.
    Using Synaptic Operations(SPOs) as the primary metric, the SNN-based controller will
    be benchmarked against traditional CNNs to validate its energy-saving potential
    in edge robotics applications.
\end{enumerate}

\subsection{Research Significance}
This research contributes to the field of Edge Artifical Intelligence(Edge AI) by
bridging the gap between biological inspiration and practical robotic engineering:

\begin{enumerate}[label=\textbf{\arabic*.}, leftmargin=*]
    \item \textbf{Practical Framework for Spike-based Robotics} \newline
    By implementing SNNs on a cost-effective, multi-processor platform(integrating STM32
    and ESP32), this work provides a scalable bileprint for the transition from theoretical
    simulation to physical deployment. It validates the robustness of event-driven vision
    in stabilizing dynamically unstable systems, offering a viable path for the development
    of long-endurance autonomous embedded devices.
    
    \item \textbf{Advancement of Software-Hardware Co-optimization} \newline
    The research demonstrates a rigorous co-optimization straegy, where the high-level C++
    SNN engine is harmomized with the low-level micro-ROS communicationlayer. 
    This approach addresses the critical constraints of Edge Intelligence,
     such as limited memory and power budgets, 
     while ensuring the millisecond-level responsiveness required for real-time control.
    
    \item \textbf{Integration of Neuromorphic Theory and Nanostructure Horizons}\newline 
    This dissertation establishes a computational foundation for the future migration of AI 
    from general-purpose CPUs to specialized neuromorphic architectures. By validating SNN 
    algorithmic efficiency on standard CMOS-based hardware, the research provides critical 
    performance benchmarks and architectural insights. These findings are essential for the 
    future design of dedicated Neuromorphic Systems-on-Chip (NSoC), facilitating the evolution 
    toward more bio-plausible and energy-efficient digital integrated circuits.
\end{enumerate}

\subsection{Thesis Structure}
The remainder of this dissertation is organized into the following seven chapters:

\begin{itemize}[leftmargin=*]

    \item \textbf{Chapter 2: Literature Review} examines the theoretical foundations of 
    Spiking Neural Networks (SNNs) alongside classical control theories. It reviews the 
    state-of-the-art in neuromorphic sensing, Real-Time Operating Systems (RTOS) for 
    robotics, and the stability challenges inherent in inverted pendulum systems.

    \item \textbf{Chapter 3: System Methodology} proposes a decoupled hierarchical architecture. 
    It details the mathematical modeling of the self-balancing robot, the design of the cascaded 
    PID control loops, and the development of a ``Retina-inspired'' visual encoding algorithm for 
    spike-based perception.

    \item \textbf{Chapter 4: Implementation and Engineering} describes the multi-processor hardware
     integration involving the STM32F103 and ESP32-S3. This chapter focuses on the embedded software 
     engineering aspects, including RTOS task prioritization, micro-ROS middleware configuration, and 
     the development of an optimized C++ SNN inference engine.

    \item \textbf{Chapter 5: Experiments and Results} presents empirical data from physical robot trials.
     It evaluates the system's performance across multiple dimensions, including balancing stability (pitch jitter),
      navigation accuracy, and computational energy efficiency quantified by Synaptic Operations (SOPs).

    \item \textbf{Chapter 6: Discussion} provides a critical analysis of the experimental findings, 
    focusing on the trade-offs between biological plausibility and real-time engineering constraints. 
    It also discusses the system's robustness in Edge AI scenarios.

    \item \textbf{Chapter 7: Conclusion} summarizes the key engineering contributions and academic insights
     of the research, offering recommendations for future advancements in neuromorphic embedded systems and 
     hardware-in-the-loop control.

\end{itemize}

\ExecuteMetaData[table]{parameters}
